% Source: tex/chapters/ch05/06_テスト技法.tex
\subsubsection{利用ベース技法}
利用ベース技法は、実際の運用や利用のされ方をモデル化し、その頻度や順序に沿ってテストケースを選ぶ。運用環境に近い条件で実行するため、受入テストや信頼性評価に向いている。

\paragraph{運用プロファイル}
運用プロファイルは、入力や利用シナリオが実運用でどの頻度で起こるかを表す。運用プロファイルに基づくテストでは、観測した結果から将来の信頼性を推定する。難しさは、現実的なプロファイルを作ることにある。

\paragraph{利用者観察ヒューリスティック}
利用者観察ヒューリスティックは、利用者がシステムやインタフェースをどれだけうまく使えるかを、統制された条件で観察する。認知的ウォークスルー、発話思考法、現場観察、質問票、面接などが使われる。
